\chapter{Nhà Em Như Vậy Thì Em Phải Sống Sao}
\markboth{Nhà Em Như Vậy Thì Em Phải Sống Sao}{How Should I Live with a Home Like Mine}

\section{Người Lớn Cãi Nhau Xong Em Vẫn Phải Đi Học}
\begin{truyen}{Người Lớn Cãi Nhau Xong Em Vẫn Phải Đi Học}{Adults Fight at Home and I Still Have to Go to School}
\chuhoa{C}{ó những học sinh đi học sau một đêm nhà đầy tiếng cãi nhau.}
\textit{(Some students come to school after a night full of adult fighting at home.)}

Sáng ra các em vẫn mặc đồng phục, vẫn chào, vẫn ngồi đúng chỗ.
\textit{(In the morning they still wear the uniform, greet people, and sit in the right seat.)}

Có em tới lớp rồi mới nói nhỏ: “Đêm qua nhà em không ngủ được.”
\textit{(Some students only whisper after arriving at school, “No one slept at home last night.”)}

Nhìn bề ngoài, chẳng có gì quá khác.
\textit{(From the outside, nothing seems very different.)}

Nhưng bên trong, em đang học trong phần dư của một đêm không yên.
\textit{(Inside, the child is learning inside the leftovers of a restless night.)}

\end{truyen}

\begin{baihoc}
Có những học sinh bước vào lớp cùng với tiếng ồn của gia đình vẫn còn ở trong người.
\textit{(Some students enter the classroom still carrying the noise of home inside them.)}
\end{baihoc}

\begin{ghinhoanh}
\textbf{Short English to remember:}

Some students come to class carrying last night's noise inside them.

\end{ghinhoanh}

\begin{tuhocnhanh}
\textbf{Gợi ý để dạy và tự nghĩ / Teaching and reflection prompts}

1. Một đêm ở nhà có thể đi theo học sinh vào lớp thế nào? \textit{How can one night at home follow a student into class?}

2. Vì sao vẻ ngoài bình thường dễ đánh lừa người lớn? \textit{Why does normal appearance mislead adults?}

3. Mình có đang nhớ học sinh mang cả gia đình vào lớp không? \textit{Do I remember that students bring their families into class with them?}
\end{tuhocnhanh}
\ngancach

\section{Nhà Em Thiếu Tiền Nên Em Thiếu Nhiều Thứ Khác}
\begin{truyen}{Nhà Em Thiếu Tiền Nên Em Thiếu Nhiều Thứ Khác}{When Home Lacks Money, Other Things Go Missing Too}
\chuhoa{N}{ghèo không chỉ thiếu tiền. Nó thường kéo theo thiếu yên tâm, thiếu thời gian của cha mẹ, thiếu không gian học, thiếu cả cảm giác mình ngang hàng với bạn bè.}
\textit{(Poverty is not only a lack of money. It often brings less security, less parental time, less study space, and less sense of equality with peers.)}

Một đứa trẻ lớn lên trong thiếu thốn phải vừa học vừa che giấu rất nhiều.
\textit{(A child growing up in hardship often studies while also hiding many things.)}

Người lớn hay khuyên cố lên.
\textit{(Adults often say try harder.)}

Nhưng có những hoàn cảnh cần được hiểu trước khi được khuyên.
\textit{(But some circumstances need understanding before advice.)}

\end{truyen}

\begin{baihoc}
Muốn công bằng với học sinh, phải nhìn thấy những bất lợi không nằm trên mặt giấy.
\textit{(If we want to be fair to students, we must see the disadvantages that do not appear on paper.)}
\end{baihoc}

\begin{ghinhoanh}
\textbf{Short English to remember:}

Fairness requires seeing burdens that are not visible on paper.

\end{ghinhoanh}

\begin{tuhocnhanh}
\textbf{Gợi ý để dạy và tự nghĩ / Teaching and reflection prompts}

1. Nghèo tác động lên việc học ngoài chuyện tiền như thế nào? \textit{How does poverty affect learning beyond money?}

2. Học sinh đang phải che giấu những gì? \textit{What may the student be hiding?}

3. Mình có đang khuyên trước khi hiểu không? \textit{Am I advising before understanding?}
\end{tuhocnhanh}
\ngancach

\section{Nếu Nhà Em Không Có Chỗ Dựa Thì Sao}
\begin{truyen}{Nếu Nhà Em Không Có Chỗ Dựa Thì Sao}{What If Home Does Not Feel Like Support}
\chuhoa{C}{ó những em không thể kể chuyện nhà mình với bạn bè vì càng kể càng thấy tủi.}
\textit{(Some students cannot talk about home with their friends because the more they tell, the more ashamed they feel.)}

Nhà với nhiều người là chỗ dựa.
\textit{(Home is support for many people.)}

Nhưng với vài đứa trẻ, nhà là chỗ các em phải tự thủ trước.
\textit{(But for some children, home is the place where they must first defend themselves.)}

Đi học với một điểm tựa lung lay là chuyện nặng hơn người ngoài nghĩ.
\textit{(Going to school with unstable support is heavier than outsiders imagine.)}

\end{truyen}

\begin{baihoc}
Không phải học sinh nào cũng có hậu phương lành để quay về sau một ngày học.
\textit{(Not every student has a healthy home base to return to after school.)}
\end{baihoc}

\begin{ghinhoanh}
\textbf{Short English to remember:}

Not every student has a safe home base.

\end{ghinhoanh}

\begin{tuhocnhanh}
\textbf{Gợi ý để dạy và tự nghĩ / Teaching and reflection prompts}

1. Vì sao thiếu điểm tựa ở nhà làm việc học nặng hơn? \textit{Why does the lack of support at home make school heavier?}

2. Học sinh thường giấu chuyện này bằng cách nào? \textit{How do students usually hide this reality?}

3. Trường học có thể bù được phần nào không? \textit{Can school partially compensate for it?}
\end{tuhocnhanh}
\ngancach

\section{Em Có Được Mong Một Cuộc Sống Khác Không}
\begin{truyen}{Em Có Được Mong Một Cuộc Sống Khác Không}{Am I Allowed to Hope for a Different Life}
\chuhoa{C}{ó những học sinh học rất bình thường, nhưng bám trường rất dai.}
\textit{(Some students perform only moderately, yet cling to school with unusual persistence.)}

Không phải vì thích học mọi môn.
\textit{(It is not because they love every subject.)}

Mà vì các em mơ một con đường khác với quãng đời người lớn quanh mình.
\textit{(It is because they dream of a different path from the adult life around them.)}

Nhiều em đến lớp không chỉ để lấy kiến thức. Các em đến để thử xem số phận có cửa rẽ nào không.
\textit{(Many students come to class not only for knowledge. They come to see whether fate has a side road.)}

\end{truyen}

\begin{baihoc}
Đôi khi điều giữ học sinh lại với trường không phải sở thích học, mà là hy vọng được sống khác.
\textit{(Sometimes what keeps students tied to school is not love of study, but the hope of a different life.)}
\end{baihoc}

\begin{ghinhoanh}
\textbf{Short English to remember:}

Some students stay in school because they hope for a different life.

\end{ghinhoanh}

\begin{tuhocnhanh}
\textbf{Gợi ý để dạy và tự nghĩ / Teaching and reflection prompts}

1. Điều gì đang giữ em này ở lại với trường? \textit{What is keeping this student in school?}

2. Hy vọng nào đang đi dưới việc học? \textit{What hope is running underneath the learning?}

3. Mình có đang tôn trọng hy vọng đó đủ không? \textit{Am I respecting that hope enough?}
\end{tuhocnhanh}
\ngancach

\section{Ở Nhà Không Ai Nghe Nên Em Nói Ồn Ở Lớp}
\begin{truyen}{Ở Nhà Không Ai Nghe Nên Em Nói Ồn Ở Lớp}{No One Hears Me at Home, So I Get Loud at School}
\chuhoa{C}{ó những học sinh ở lớp nói rất chen, rất nhiều, rất thích giành lời.}
\textit{(Some students talk over others in class, speak too much, and constantly try to take the floor.)}

Người lớn nhìn vào thấy phiền.
\textit{(Adults naturally find it annoying.)}

Nhưng có em nói thật: “Ở nhà em nói không ai nghe.”
\textit{(But some students say honestly, “No one listens to me at home.”)}

Có những cái ồn ở trường là phần dội lại của một sự im ở nhà.
\textit{(Some noise at school is the echo of being unheard at home.)}

\end{truyen}

\begin{baihoc}
Muốn hiểu một học sinh nói ồn, đôi khi phải nhìn vào chỗ em không được nghe ở nơi khác.
\textit{(To understand a loud student, adults sometimes must look at where the child is not heard elsewhere.)}
\end{baihoc}

\begin{ghinhoanh}
\textbf{Short English to remember:}

Some noise at school is the echo of being unheard at home.

\end{ghinhoanh}

\begin{tuhocnhanh}
\textbf{Gợi ý để dạy và tự nghĩ / Teaching and reflection prompts}

1. Vì sao bị không được nghe ở nhà có thể thành ồn ở lớp? \textit{Why can being unheard at home become noise in class?}

2. Mình đang chỉ thấy phần phiền hay thấy cả nguồn gốc? \textit{Am I seeing only the disturbance or also the source?}

3. Người dạy có thể cho em một chỗ nói khác đi bằng cách nào? \textit{How can a teacher offer the student another way to be heard?}
\end{tuhocnhanh}
\ngancach

\section{Nhà Em Có Chuyện Nên Em Học Không Nổi}
\begin{truyen}{Nhà Em Có Chuyện Nên Em Học Không Nổi}{There Is Trouble at Home, So I Cannot Study Properly}
\chuhoa{C}{ó học sinh bị tụt rất nhanh trong vài tuần. Hỏi kỹ mới biết nhà đang có biến cố lớn.}
\textit{(Some students decline very quickly over a few weeks. A closer conversation reveals a major problem at home.)}

Người ngoài chỉ thấy em sao lơ là vậy.
\textit{(Outsiders only ask why the student has become so careless.)}

Nhưng đầu óc một đứa trẻ không thể cùng lúc ôm hết chuyện nhà và học như không có gì.
\textit{(But a child's mind cannot carry a family crisis and still study as if nothing is happening.)}

Có giai đoạn học lực tụt không phải vì kém đi, mà vì đời sống đang dồn sang hướng khác.
\textit{(There are times academic decline is not a drop in ability, but life pressing from another side.)}

\end{truyen}

\begin{baihoc}
Khi nhà có chuyện, bài vở thường là nơi lộ ra trước tiên rằng học sinh đang quá tải.
\textit{(When home is in crisis, schoolwork is often the first place where overload becomes visible.)}
\end{baihoc}

\begin{ghinhoanh}
\textbf{Short English to remember:}

Schoolwork often reveals first when family life is too heavy.

\end{ghinhoanh}

\begin{tuhocnhanh}
\textbf{Gợi ý để dạy và tự nghĩ / Teaching and reflection prompts}

1. Vì sao chuyện nhà thường lộ ra qua việc học? \textit{Why does family trouble often show up in school performance?}

2. Người lớn dễ kết luận sai ở điểm nào? \textit{Where do adults most easily misread the situation?}

3. Mình có đang hỏi đủ phần đời sống ngoài lớp chưa? \textit{Am I asking enough about life outside class?}
\end{tuhocnhanh}
\ngancach

\section{Em Không Dám Mời Bạn Về Nhà}
\begin{truyen}{Em Không Dám Mời Bạn Về Nhà}{I Do Not Dare Invite Friends Home}
\chuhoa{C}{ó học sinh luôn kiếm cớ khi bạn rủ về chơi qua lại.}
\textit{(Some students always find excuses when friends suggest visiting each other's homes.)}

Không phải vì các em không quý bạn.
\textit{(It is not because they do not care about their friends.)}

Các em sợ người khác thấy nhà mình chật, nghèo, ồn, hoặc quá buồn.
\textit{(They fear others seeing their house as cramped, poor, noisy, or too sad.)}

Tuổi học trò có những xấu hổ rất nhỏ mà đè rất nặng.
\textit{(School age carries small forms of shame that weigh heavily.)}

\end{truyen}

\begin{baihoc}
Nhiều học sinh phải giấu hoàn cảnh sống để giữ lòng tự trọng giữa bạn bè.
\textit{(Many students hide their living conditions to protect dignity among peers.)}
\end{baihoc}

\begin{ghinhoanh}
\textbf{Short English to remember:}

Some students hide home to protect their dignity.

\end{ghinhoanh}

\begin{tuhocnhanh}
\textbf{Gợi ý để dạy và tự nghĩ / Teaching and reflection prompts}

1. Vì sao việc mời bạn về nhà lại nặng với một số học sinh? \textit{Why is inviting friends home so heavy for some students?}

2. Lòng tự trọng ở đây đang được giữ bằng cách nào? \textit{How is dignity being protected here?}

3. Mình có đủ tinh tế với những hoàn cảnh kiểu này không? \textit{Am I sensitive enough to this kind of situation?}
\end{tuhocnhanh}
\ngancach

\section{Nhà Em Không Có Ai Chỉ Bài}
\begin{truyen}{Nhà Em Không Có Ai Chỉ Bài}{There Is No One at Home to Help Me Study}
\chuhoa{C}{ó em về nhà là tự xoay hết: bài vở, ăn uống, em nhỏ, việc lặt vặt.}
\textit{(Some students go home and must handle everything alone: homework, meals, younger siblings, and daily chores.)}

Người lớn ở trường bảo về nhà hỏi thêm cha mẹ.
\textit{(Adults at school tell them to ask their parents for help.)}

Nhưng có những nhà không có ai đủ thời gian, đủ sức, hay đủ học để chỉ thêm.
\textit{(But some homes do not have anyone with enough time, energy, or education to help.)}

Khi ấy, việc học của trẻ không chỉ là học, mà là tự gắng một mình rất sớm.
\textit{(Then learning becomes not just learning, but an early form of surviving alone.)}

\end{truyen}

\begin{baihoc}
Không phải học sinh nào cũng có một hậu phương học tập ở nhà.
\textit{(Not every student has an academic support system at home.)}
\end{baihoc}

\begin{ghinhoanh}
\textbf{Short English to remember:}

Not every student has study support waiting at home.

\end{ghinhoanh}

\begin{tuhocnhanh}
\textbf{Gợi ý để dạy và tự nghĩ / Teaching and reflection prompts}

1. Vì sao lời “về nhà hỏi thêm” không phải lúc nào cũng khả thi? \textit{Why is “ask for help at home” not always realistic?}

2. Học sinh đang phải tự xoay những gì? \textit{What is the student managing alone?}

3. Trường học có thể bù phần thiếu này ra sao? \textit{How can school help fill this gap?}
\end{tuhocnhanh}
\ngancach

\section{Có Những Đứa Trẻ Đi Học Như Một Cách Tạm Trốn Nhà}
\begin{truyen}{Có Những Đứa Trẻ Đi Học Như Một Cách Tạm Trốn Nhà}{Some Children Come to School as a Temporary Escape from Home}
\chuhoa{C}{ó học sinh đi học rất đều không hẳn vì mê học.}
\textit{(Some students come to school very regularly not because they love lessons.)}

Các em đi vì ở trường dễ thở hơn ở nhà.
\textit{(They come because school feels easier to breathe in than home.)}

Điều này vừa buồn vừa quan trọng.
\textit{(This is both sad and important.)}

Nó nhắc rằng với một số đứa trẻ, lớp học không chỉ là nơi học mà còn là nơi tạm trú tinh thần.
\textit{(It reminds us that for some children, the classroom is not only a place of learning but also temporary emotional shelter.)}

\end{truyen}

\begin{baihoc}
Với vài học sinh, trường học có giá trị lớn vì nó là nơi các em đỡ ngột hơn phần đời còn lại.
\textit{(For some students, school matters because it is the place where life feels less suffocating.)}
\end{baihoc}

\begin{ghinhoanh}
\textbf{Short English to remember:}

For some children, school is also temporary shelter.

\end{ghinhoanh}

\begin{tuhocnhanh}
\textbf{Gợi ý để dạy và tự nghĩ / Teaching and reflection prompts}

1. Vì sao trường học lại thành chỗ tạm trú với vài em? \textit{Why can school become temporary shelter for some students?}

2. Điều này đặt thêm trách nhiệm gì lên lớp học? \textit{What added responsibility does this place on the classroom?}

3. Mình có đang hiểu hết ý nghĩa của việc các em đi học đều không? \textit{Am I fully understanding what regular attendance may mean?}
\end{tuhocnhanh}
\ngancach

\section{Em Không Chọn Được Nhà Mình, Nhưng Em Muốn Chọn Được Đời Mình}
\begin{truyen}{Em Không Chọn Được Nhà Mình, Nhưng Em Muốn Chọn Được Đời Mình}{I Cannot Choose My Home, but I Want to Choose My Life}
\chuhoa{C}{ó em nói một câu rất chín: “Nhà em vậy rồi. Em chỉ mong đời em đừng y như vậy.”}
\textit{(One student said something very mature: “My home is already like this. I only hope my life won't stay exactly like this.”)}

Trong một câu ngắn có đủ buồn, đủ tỉnh, và đủ hy vọng.
\textit{(That short sentence holds sadness, clarity, and hope all at once.)}

Giáo dục nhiều khi không cứu được hoàn cảnh hiện tại ngay.
\textit{(Education often cannot rescue the present circumstance immediately.)}

Nhưng nó có thể giúp học sinh tin rằng mình còn có quyền chọn phần đường phía trước.
\textit{(But it can help students believe they still have some power over the road ahead.)}

\end{truyen}

\begin{baihoc}
Một trong những món quà lớn của giáo dục là trả lại cho học sinh niềm tin rằng cuộc đời mình chưa bị khóa chặt bởi hoàn cảnh xuất thân.
\textit{(One of education's great gifts is restoring the belief that life is not locked forever by starting circumstances.)}
\end{baihoc}

\begin{ghinhoanh}
\textbf{Short English to remember:}

Education can restore the belief that life is still open ahead.

\end{ghinhoanh}

\begin{tuhocnhanh}
\textbf{Gợi ý để dạy và tự nghĩ / Teaching and reflection prompts}

1. Vì sao hy vọng này lại quan trọng sống còn với một số học sinh? \textit{Why is this hope vital for some students?}

2. Giáo dục có thể cho các em điều gì dù chưa đổi được hoàn cảnh ngay? \textit{What can education give even before circumstances change?}

3. Mình có đang giúp học sinh tin vào phần đời phía trước không? \textit{Am I helping students believe in the road ahead?}
\end{tuhocnhanh}
